%% Springer Nature LaTeX template — 测试夹具（最小可用骨架，非官方原件）
%%
%% 这套模板的元信息写法跟 Elsevier CAS 完全不同，是回归测试要钉住的形态：
%%   * `\author*[1,2]{}` / `\affil*[1]{}` **带星号**（通讯作者 / 第一单位）
%%   * 姓名拆成 `\fnm{名} \sur{姓}`
%%   * 单位拆成 `\orgdiv{}` / `\orgname{}` / `\orgaddress{}`
%%   * 摘要与关键词都是**命令式**（`\abstract{}` / `\keywords{}`），不是环境
%%   * `.bst` 放在 `bst/` 子目录里
\documentclass[pdflatex,sn-mathphys-num]{sn-jnl}

\begin{document}

\title[Article Title]{Article Title}

%%=============================================================%%
%% GivenName	-> \fnm{Joergen W.}
%% FamilyName	-> \sur{Ploeg}
%%=============================================================%%

\author*[1,2]{\fnm{First} \sur{Author}}\email{iauthor@gmail.com}

\author[2,3]{\fnm{Second} \sur{Author}}\email{iiauthor@gmail.com}

\affil*[1]{\orgdiv{Department}, \orgname{Organization}, \orgaddress{\street{Street}, \city{City}, \postcode{100190}, \state{State}, \country{Country}}}

\affil[2]{\orgdiv{Department}, \orgname{Organization}, \orgaddress{\street{Street}, \city{City}, \postcode{10587}, \state{State}, \country{Country}}}

\abstract{The abstract serves both as a general introduction to the topic and as a brief, non-technical summary of the main results and their implications. Authors are advised to check the author instructions for the journal they are submitting to for word limits.}

\keywords{keyword1, Keyword2, Keyword3, Keyword4}

\maketitle

\section{Introduction}\label{sec1}

Template sample body text that must be replaced by the source manuscript body.

\section{Results}\label{sec2}

More template sample body text.

\bibliography{sn-bibliography}

\end{document}
